\documentclass[10pt]{scrartcl}
\usepackage[mdthm,colorsec,noauthor,noasy]{evan}
\usepackage[letterpaper,margin=0.75in,headsep=0.25in]{geometry}
\usepackage{needspace}
\setlist{itemsep=2pt,topsep=2pt,parsep=0pt}
\setlength{\parindent}{0pt}
\setlength{\emergencystretch}{3em}
\DeclareMathOperator{\nul}{null}
\DeclareMathOperator{\range}{range}
\DeclareMathOperator{\spn}{span}
\DeclareMathOperator{\tr}{tr}
\newcommand{\LL}{\mathcal L}
\newcommand{\PP}{\mathcal P}
\newcommand{\MM}{\mathcal M}
\newcommand{\ip}[2]{\langle #1,#2\rangle}
\newcommand{\idea}[1]{{\color{OliveGreen!75!black}\textbf{Idea.}} \textit{#1}\par}
\newcommand{\Q}[2]{\par\needspace{5\baselineskip}\medskip\noindent\textbf{#1. #2}\par}
\newenvironment{prob}{\par\medskip\noindent\begin{minipage}{\linewidth}\setlength{\parskip}{3pt}}{\end{minipage}\par}
\title{Math 110 Review}
\subtitle{Solutions}
\author{}
\date{}
\ihead{}
\begin{document}
\maketitle
\setcounter{section}{1}

\section{Problem set}

\needspace{14\baselineskip}
\subsection{Minimal polynomial (HW7)}

\begin{prob}
\textbf{HW7 P1. }\par 
\idea{$A^2=-I$ and $(B-I)^2=0$ give low-degree polynomials that kill $A$ and $B$; then rule out degree $1$.}
\begin{enumerate}[(a)]
  \ii $A^2=-I$, so $m_A\mid x^2+1$. $A$ is not a multiple of $I$ ($Ae_1=e_2$ is not a multiple of $e_1$), so $\deg m_A\ne1$. Hence $m_A=x^2+1$ over $\RR$. Over $\CC$ the definition (first $k$ with $A^k\in\spn(I,\dots,A^{k-1})$) gives the same polynomial: $x^2+1=(x-i)(x+i)$.
  \ii $(B-I)^2=\begin{psmallmatrix}0&1\\0&0\end{psmallmatrix}^2=0$ and $B\ne I$, so $m_B=(x-1)^2$.
  \ii $A$ over $\RR$: $x^2+1$ has no real roots, so it is not a product of real linear factors: not diagonalizable. Over $\CC$: distinct roots $\pm i$, diagonalizable. $B$: repeated root, not diagonalizable. This matches Homework 6.
\end{enumerate}
\end{prob}

\begin{prob}
\textbf{HW7 P2. }\par 
\idea{$D$ lowers degree, so a power of it is $0$; find the first one.}
(a) $D^4=0$ on $\PP_3(\RR)$, while $D^3x^3=6\ne0$. So $m_D\mid x^4$ and is not $x^k$ for $k\le3$: $m_D=x^4$. (b) Repeated root: not diagonalizable. (Also $E(0,D)=$ constants, of dimension $1<4$.)
\end{prob}

\begin{prob}
\textbf{HW7 P3. }\par 
\begin{enumerate}[(a)]
  \ii \idea{If $a_0=0$, then $x$ divides $m_T$; since $T$ is invertible, the factor $T$ can be cancelled, leaving a smaller polynomial that kills $T$.}
  If $a_0=0$, then $m_T=xq$ with $\deg q=d-1$, and $0=m_T(T)=Tq(T)$. Multiply by $T^{-1}$: $q(T)=0$, contradicting minimality of $m_T$ (if $d=1$, $q=1$ and $q(T)=I\ne0$).
  \ii \idea{Move the constant term to one side and factor $T$ out of the rest.}
  From $m_T(T)=0$: $T\big(T^{d-1}+a_{d-1}T^{d-2}+\dots+a_1I\big)=-a_0I$. So $T^{-1}=-\frac1{a_0}\big(T^{d-1}+a_{d-1}T^{d-2}+\dots+a_1I\big)$.
  \ii \idea{Apply (b), then factor $m_T$ to read off the eigenvalues.}
  $a_0=4$, so $T^{-1}=-\frac14(T^3+T^2-T-5I)$. $m_T=(x-1)^2(x^2+3x+4)$ (check: $1$ is a root of $m_T$ and of $m_T'=4x^3+3x^2-2x-5$). The eigenvalues are $1$ and $\frac{-3\pm i\sqrt7}2$. The root $1$ is repeated, so $T$ is not diagonalizable.
\end{enumerate}
\end{prob}

\begin{prob}
\textbf{HW7 P4. }\par 
\idea{Substituting $x\mapsto x+\lambda$ (or $x\mapsto x/c$) turns polynomials that kill $T$ into polynomials that kill $T-\lambda I$ (or $cT$), and back. So the two minimal polynomials divide each other's transforms and have the same degree.}
\begin{enumerate}[(a)]
  \ii $p(x)=m_T(x+\lambda)$ is monic of degree $d=\deg m_T$ and $p(T-\lambda I)=m_T(T)=0$, so $m_{T-\lambda I}\mid p$ and $\deg m_{T-\lambda I}\le d$. Likewise $r(x)=m_{T-\lambda I}(x-\lambda)$ is monic with $r(T)=m_{T-\lambda I}(T-\lambda I)=0$, so $d\le\deg m_{T-\lambda I}$. A monic divisor of $p$ of the same degree is $p$.
  \ii $p(x)=c^dm_T(x/c)$ is monic of degree $d$ (leading coefficient $c^d\cdot c^{-d}$) and $p(cT)=c^dm_T(T)=0$, so $m_{cT}\mid p$. The same argument with $cT$ and $1/c$ gives $d\le\deg m_{cT}$. So $m_{cT}=p$.
\end{enumerate}
\end{prob}

\begin{prob}
\textbf{HW7 P5. }\par 
\idea{The roots of $m_A$ are the eigenvalues and $\deg m_A\le2$, which leaves two candidates.}
$\deg m_A\le2$ (Cayley--Hamilton) and its only root is $\lambda$, and over $\CC$ it splits, so $m_A=x-\lambda$ or $(x-\lambda)^2$. $m_A=x-\lambda$ exactly when $A-\lambda I=0$. $A$ is diagonalizable iff $m_A$ has no repeated root iff $m_A=x-\lambda$ iff $A=\lambda I$.
\end{prob}

\begin{prob}
\textbf{HW7 P6. }\par 
\idea{Triangular, so the eigenvalues are on the diagonal; list the candidate minimal polynomials and test them.}
Eigenvalues $2,3$. $E(2)=\nul\begin{psmallmatrix}0&1&0\\0&0&0\\0&0&1\end{psmallmatrix}=\spn(e_1)$, $E(3)=\nul\begin{psmallmatrix}-1&1&0\\0&-1&0\\0&0&0\end{psmallmatrix}=\spn(e_3)$. Candidates (roots exactly $2,3$, degree $\le3$): $(x-2)(x-3)$, $(x-2)^2(x-3)$, $(x-2)(x-3)^2$. $(C-2I)(C-3I)$ and $(C-2I)(C-3I)^2$ have $(1,2)$ entry $\mp1\ne0$; $(C-2I)^2=\operatorname{diag}(0,0,1)$ and $(C-2I)^2(C-3I)=0$. So $m_C=(x-2)^2(x-3)$: not diagonalizable ($\dim E(2)+\dim E(3)=2<3$).
\end{prob}

\begin{prob}
\textbf{HW7 P7. }\par 
\idea{Lagrange interpolation: a product of factors $(x-\lambda_k)$ vanishes at every $\lambda_k$ but $\lambda_j$; divide by its value at $\lambda_j$.}
$p_j(x)=\prod_{k\ne j}\frac{x-\lambda_k}{\lambda_j-\lambda_k}$. Since $m_T$ has distinct roots, $T$ is diagonalizable and $V=\bigoplus_kE(\lambda_k,T)$. For $v\in E(\lambda_k,T)$, $p_j(T)v=p_j(\lambda_k)v$, which is $v$ if $k=j$ and $0$ otherwise. So for $v=\sum_kv_k$, $p_j(T)v=v_j$: the projection onto $E(\lambda_j,T)$ along $\bigoplus_{k\ne j}E(\lambda_k,T)$.
\end{prob}

\begin{prob}
\textbf{M1. }\par 
\idea{Find a low-degree polynomial relation, then rule out its proper divisors. In (c), $B$ is already in Jordan form, so read $m_B$ off the blocks: the largest block for each eigenvalue gives its exponent.}
\begin{enumerate}[(a)]
  \ii Every entry of $A^2$ is $3$, so $A^2=3A$ and $m_A\mid x(x-3)$. $A\ne0$ and $A\ne3I$, so $m_A=x(x-3)$. Distinct roots: diagonalizable.
  \ii $Pe_1=e_2$, $Pe_2=e_3$, $Pe_3=e_1$, so $P^3=I$ and $m_P\mid x^3-1$. Since $e_1,Pe_1,P^2e_1$ is a basis, no nonzero polynomial of degree $\le2$ kills $e_1$, so $m_P=x^3-1=(x-1)(x^2+x+1)$. Over $\CC$ the roots $1,\omega,\bar\omega$ are distinct: diagonalizable. Over $\RR$, $x^2+x+1$ has no real root: not diagonalizable.
  \ii $B=J_2(3)\oplus J_1(3)\oplus J_1(5)$, so $m_B=\operatorname{lcm}\big((x-3)^2,x-3,x-5\big)=(x-3)^2(x-5)$. Repeated root: not diagonalizable.
\end{enumerate}
\end{prob}

\begin{prob}
\textbf{M2. }\par 
\begin{enumerate}[(a)]
  \ii \idea{Test the basis vectors: $x^k$ is an eigenvector.}
  $T(x^k)=kx^k$, so $1,x,\dots,x^n$ is an eigenbasis with distinct eigenvalues $0,1,\dots,n$. $T$ is diagonalizable and $m_T=x(x-1)\cdots(x-n)$.
  \ii \idea{$S-I$ is a ``difference operator'' that lowers degree, like differentiation; so it is nilpotent, and the question is the first power that kills everything.}
  $(S-I)p=p(x+1)-p(x)$: the leading terms cancel, so $\deg$ drops by at least $1$ and constants go to $0$. Hence $(S-I)^{n+1}=0$, so $m_S\mid(x-1)^{n+1}$. Also $(S-I)x^k=(x+1)^k-x^k=kx^{k-1}+(\text{lower})$, so $(S-I)^nx^n=n!\ne0$. So $m_S=(x-1)^{n+1}$, a repeated root: not diagonalizable.
\end{enumerate}
\end{prob}

\begin{prob}
\textbf{M3. }\par 
\idea{$T^3=T$ says $x^3-x$ kills $T$; factor it.}
$m_T\mid x^3-x=x(x-1)(x+1)$, so $m_T$ is a product of distinct real linear factors and $T$ is diagonalizable. Its eigenvalues are roots of $m_T$, hence in $\{-1,0,1\}$.
\end{prob}

\begin{prob}
\textbf{M4. }\par 
\idea{A monic polynomial of degree $1$ is $x-c$, and $m_T=x-c$ says $T-cI=0$.}
If $T=cI$, then $x-c$ kills $T$ and $\deg m_T\ge1$ (as $V\ne0$), so $m_T=x-c$. Conversely, if $m_T=x-c$, then $T-cI=m_T(T)=0$.
\end{prob}

\begin{prob}
\textbf{M5. }\par 
\idea{Substitute $x\mapsto1/x$ in $m_T(T)=0$ and clear denominators by multiplying by $T^{-d}$.}
$q(x)=\frac1{a_0}x^dm_T(1/x)=\frac1{a_0}\big(1+a_{d-1}x+\dots+a_0x^d\big)$ is monic of degree $d$ ($a_0\ne0$ by HW7 P3). $q(T^{-1})=\frac1{a_0}T^{-d}m_T(T)=0$, so $m_{T^{-1}}\mid q$ and $\deg m_{T^{-1}}\le d$. Applying this to $T^{-1}$ in place of $T$ gives $d\le\deg m_{T^{-1}}$. So $m_{T^{-1}}=q$.
\end{prob}

\begin{prob}
\textbf{M6. }\par 
\idea{Let $r_i$ be the step where the chain $\nul(T-\lambda_iI)^k$ reaches $G(\lambda_i,T)$. $\prod(x-\lambda_j)^{r_j}$ kills every $G$, giving $k_i\le r_i$. For $\ge$: if $k_i<r_i$, some vector of $G(\lambda_i,T)$ survives $(T-\lambda_iI)^{k_i}$, and the other factors of $m_T$ are injective on $G(\lambda_i,T)$, so $m_T(T)$ would not kill it.}
Let $r_i$ be the smallest $k$ with $\nul(T-\lambda_iI)^k=G(\lambda_i,T)$.
\begin{enumerate}[1.]
  \ii $q=\prod_j(x-\lambda_j)^{r_j}$ kills each $G(\lambda_j,T)$ (factors commute), and these span $V$ (Prop.~73; by (T9) the $\lambda_j$ are all the eigenvalues of $T$). So $m_T\mid q$ and $k_i\le r_i$.
  \ii For $j\ne i$, $T-\lambda_jI$ is injective on $G(\lambda_i,T)$: a vector it kills is an eigenvector for $\lambda_j$ lying in $G(\lambda_i,T)$, and generalized eigenvectors for distinct eigenvalues are independent (Prop.~72), so it is $0$.
  \ii If $k_i<r_i$, pick $v\in G(\lambda_i,T)$ with $w=(T-\lambda_iI)^{k_i}v\ne0$; $w\in G(\lambda_i,T)$ (invariant). Then $0=m_T(T)v=\prod_{j\ne i}(T-\lambda_jI)^{k_j}w\ne0$ by step 2. Contradiction, so $k_i=r_i$.
\end{enumerate}
In a Jordan basis, $\nul(T-\lambda_iI)^k=G(\lambda_i,T)$ exactly when $k$ is at least the largest block size for $\lambda_i$.
\end{prob}

\begin{prob}
\textbf{M7. }\par 
\idea{Block-count table: differences of the $d_k$ count blocks of size $\ge k$.}
$d_0,\dots,d_5=0,3,5,6,7,7$ ($d_5=7$ since $\dim V=7$). Blocks of size $\ge1$: $3$; $\ge2$: $2$; $\ge3$: $1$; $\ge4$: $1$; $\ge5$: $0$. Sizes $4,2,1$ (sum $7$). $T\sim J_4(4)\oplus J_2(4)\oplus J_1(4)$, $m_T=(x-4)^4$, $\chi_T=(x-4)^7$.
\end{prob}

\begin{prob}
\textbf{M8. }\par 
\idea{Translate: $\dim E(1)$ is the number of blocks for $1$; the exponent in $m_T$ is the largest block; the total dimension fixes the rest.}
Eigenvalue $1$: one block, of size $2$. Eigenvalue $-2$: blocks of size $1$, total $4-2=2$. So $T\sim J_2(1)\oplus J_1(-2)\oplus J_1(-2)$, $\chi_T=(x-1)^2(x+2)^2$, $\tr T=1+1-2-2=-2$.
\end{prob}

\begin{prob}
\textbf{M9. }\par 
\idea{$m_A(A)=0$ is a rule for rewriting $A^2$ in terms of $A$ and $I$; use it to reduce any power, and solve it for $A^{-1}$.}
$A$ is not scalar, so $m_A=\chi_A=x^2-(\tr A)x+\det A=x^2-5x-2$. From $A^2-5A=2I$: $A(A-5I)=2I$, so $A^{-1}=\frac12(A-5I)=\begin{psmallmatrix}-2&1\\3/2&-1/2\end{psmallmatrix}$. $A^3=A(5A+2I)=5A^2+2A=5(5A+2I)+2A=27A+10I$.
\end{prob}

\needspace{14\baselineskip}
\subsection{Generalized eigenvectors and Jordan form (HW8)}

\begin{prob}
\textbf{HW8 P1. }\par 
\begin{enumerate}[(a)]
  \ii \idea{Differentiation lowers degree, so enough derivatives kill any polynomial; the index is the first power that kills $x^n$, the slowest vector to die.}$D$ lowers degree by $1$ and kills constants, so $D^{n+1}p=0$ for every $p\in\PP_n(\CC)$, while $D^nx^n=n!\ne0$. So $D$ is nilpotent of index $n+1$.
  \ii $E(0,D)=\nul D=$ constants, so the geometric multiplicity is $1$. $G(0,D)=\PP_n(\CC)$ since $D$ is nilpotent, so the algebraic multiplicity is $n+1$.
  \ii \idea{One Jordan block means one chain. Start from the vector that takes longest to die ($x^n$) and differentiate repeatedly; dividing by factorials makes each step land on the next basis vector with coefficient $1$.}$f_k=x^k/k!$ for $k=0,\dots,n$: $Df_0=0$ and $Df_k=x^{k-1}/(k-1)!=f_{k-1}$. In the basis $f_0,f_1,\dots,f_n$, column $k$ of the matrix has a single $1$ in row $k-1$: the matrix is the single block $J_{n+1}(0)$.
\end{enumerate}
\end{prob}

\begin{prob}
\textbf{HW8 P2. }\par 
\begin{enumerate}[(a)]
  \ii \idea{Translate the data into blocks: geometric multiplicity $=$ number of blocks, algebraic $=$ total size, largest block $=$ exponent in $m$.}Number of blocks $=$ geometric multiplicity, total size $=$ algebraic multiplicity: $J_3(2)\oplus J_1(-1)\oplus J_1(-1)$,
  \[ \begin{psmallmatrix}2&1&0&0&0\\0&2&1&0&0\\0&0&2&0&0\\0&0&0&-1&0\\0&0&0&0&-1\end{psmallmatrix},\qquad m=(x-2)^3(x+1). \]
  \ii Two blocks for $2$ with total size $3$: $J_2(2)\oplus J_1(2)\oplus J_1(-1)\oplus J_1(-1)$, i.e.\ the matrix in (a) with the $(2,3)$ entry changed to $0$. The largest block for $2$ has size $2$, so $m=(x-2)^2(x+1)$.
\end{enumerate}
\end{prob}

\begin{prob}
\textbf{HW8 P3. }\par 
\begin{enumerate}[(a)]
  \ii $N=J-I$ has $Ne_1=0$, $Ne_2=e_1$, $Ne_3=e_2$. So $\nul N=\spn(e_1)$; $N^2e_3=e_1$ and $N^2e_1=N^2e_2=0$, so $\nul N^2=\spn(e_1,e_2)$; $N^3=0$, so $\nul N^3=\FF^3$.
  \ii $G(1,J)=\nul(J-I)^3=\FF^3$. Geometric multiplicity $\dim\nul(J-I)=1$; algebraic multiplicity $\dim G(1,J)=3$.
  \ii \idea{Split $J$ into a scalar part and a nilpotent part that commute, then use the binomial theorem; it stops at $N^2$ because $N^3=0$.}$J=I+N$ with $IN=NI$ and $N^3=0$, so $J^m=I+mN+\binom m2N^2=\begin{psmallmatrix}1&m&\binom m2\\0&1&m\\0&0&1\end{psmallmatrix}$ (with $\binom m2=\frac{m(m-1)}2$; for $m=1$ this is $J$).
\end{enumerate}
\end{prob}

\begin{prob}
\textbf{HW8 P4. }\par 
\begin{enumerate}[(a)]
  \ii \idea{Null spaces of a block diagonal matrix add block by block, and for one block $J_j(0)^k$ is a shift by $k$ places.}For one block $J_j(0)$, $\dim\nul J_j(0)^k=\min(k,j)$. Add over the blocks of sizes $3$ and $2$: $k=0$: $0$; $k=1$: $1+1=2$; $k=2$: $2+2=4$; $k=3$: $3+2=5$; $k=4$: $3+2=5$.
  \ii \idea{Differences of the $d_k$ count blocks of size $\ge k$; the total dimension pins down the largest size.}The number of blocks of size $\ge k$ is $\dim\nul N^k-\dim\nul N^{k-1}$: $2-0=2$ blocks of size $\ge1$, $4-2=2$ of size $\ge2$, $5-4=1$ of size $\ge3$. So there are two blocks, both of size at least $2$, exactly one of size at least $3$. The three numbers alone do not fix the larger size ($J_4(0)\oplus J_2(0)$ on $\CC^6$ also gives $2,4,5$), so we use that $N$ acts on the same $5$-dimensional space as in (a): then $\nul N^3=V$ and $N^3=0$, so no block has size above $3$. Hence the sizes are $3$ and $2$. (Equivalently: the sizes are $2$ and $s\ge3$, and $2+s=\dim V=5$.)
\end{enumerate}
\end{prob}

\begin{prob}
\textbf{HW8 P5. }\par 
\idea{For each eigenvalue, make the table of $d_k$ and take differences (Prop.~77).}Eigenvalue $1$: $d_1,d_2,d_3=1,2,2$, so $1$ block of size $\ge1$, $1$ of size $\ge2$, none of size $\ge3$: $J_2(1)$. Eigenvalue $2$: $d_1,\dots,d_4=2,3,4,4$, so $2$ blocks of size $\ge1$, $1$ of size $\ge2$, $1$ of size $\ge3$, none of size $\ge4$: $J_3(2)$ and $J_1(2)$. Check: $2+3+1=6$. So $T\sim J_2(1)\oplus J_3(2)\oplus J_1(2)$.
\end{prob}

\begin{prob}
\textbf{HW8 P6. }\par 
(a) $\chi=(x-4)^5(x+1)$. (b) $m=(x-4)^3(x+1)$ (the largest block for $4$ has size $3$). (c) Eigenvalue $4$: geometric $2$ (two blocks), algebraic $5$; eigenvalue $-1$: geometric $1$, algebraic $1$. (d) Not diagonalizable: $m$ has the repeated root $4$ (equivalently, $2\ne5$).
\end{prob}

\begin{prob}
\textbf{HW8 P7. }\par 
\idea{Everything about a block diagonal matrix splits into blocks, so count block by block; for one block the null space of $J_j(0)^k$ is easy because $J_j(0)^k$ is a shift.}Choose a basis in which $N$ is in Jordan form with blocks $J_{j_1}(0),\dots,J_{j_r}(0)$. $N^k$ is block diagonal with blocks $J_{j_i}(0)^k$, so $\dim\nul N^k=\sum_i\dim\nul J_{j_i}(0)^k$. On one block with basis $e_1,\dots,e_j$, $J_j(0)^ke_l=e_{l-k}$ for $l>k$ and $0$ for $l\le k$; so $\nul J_j(0)^k=\spn(e_1,\dots,e_{\min(k,j)})$ and $\dim\nul J_j(0)^k=\min(k,j)$. Then $\min(k,j)-\min(k-1,j)$ is $1$ if $j\ge k$ and $0$ if $j<k$. Summing over the blocks, $\dim\nul N^k-\dim\nul N^{k-1}$ counts the blocks with $j_i\ge k$.
\end{prob}

\begin{prob}
\textbf{G1. }\par 
\idea{Get the eigenvalues from a characteristic polynomial, then make the $d_k$ table, then build chains from the top.}
The upper-left block has trace $4$ and determinant $4$, so $\chi_A=(x-2)^2\cdot(x-2)=(x-2)^3$: only eigenvalue $2$. $A-2I=\begin{psmallmatrix}-1&1&0\\-1&1&0\\0&0&0\end{psmallmatrix}$ has rank $1$, so $E(2,A)=\{y=x\}=\spn\big((1,1,0),(0,0,1)\big)$ and $d_1=2$. $(A-2I)^2=0$, so $d_2=3$ and $G(2,A)=\CC^3$. Blocks: $2$ of size $\ge1$, $1$ of size $\ge2$: $J_2(2)\oplus J_1(2)$, $m_A=(x-2)^2$. Jordan basis: top $e_1\notin\nul(A-2I)$, bottom $(A-2I)e_1=(-1,-1,0)$; then the eigenvector $e_3$, independent of that bottom. In the basis $(-1,-1,0),e_1,e_3$ the matrix is $J_2(2)\oplus J_1(2)$ (check: $Ae_1=(1,-1,0)=2e_1+(-1,-1,0)$).
\end{prob}

\begin{prob}
\textbf{G2. }\par 
\idea{Same as for eigenspaces: $S$ commutes with every polynomial in $T$, so it preserves null spaces of polynomials in $T$.}
$G(\lambda,T)=\nul(T-\lambda I)^n$. If $(T-\lambda I)^nv=0$, then $(T-\lambda I)^nSv=S(T-\lambda I)^nv=0$.
\end{prob}

\begin{prob}
\textbf{G3. }\par 
\begin{enumerate}[(a)]
  \ii \idea{$(\lambda I+N)v=0$ says $v$ is an eigenvector of $N$ for $-\lambda$, and $0$ is the only eigenvalue of $N$.}
  If $\lambda\ne0$ and $(\lambda I+N)v=0$, then $Nv=-\lambda v$; since $-\lambda$ is not an eigenvalue of $N$, $v=0$. So $\lambda I+N$ is injective, hence invertible. If $\lambda=0$, $N$ is not injective ($V\ne0$ and $N^k=0$). (Explicitly, $(\lambda I+N)^{-1}=\lambda^{-1}\sum_{j=0}^{n-1}(-N/\lambda)^j$.)
  \ii \idea{With one eigenvalue there is one generalized eigenspace, and it is everything.}
  By Prop.~73, $V=G(\lambda,T)=\nul(T-\lambda I)^n$, so $(T-\lambda I)^n=0$.
\end{enumerate}
\end{prob}

\begin{prob}
\textbf{G4. }\par 
\idea{Block sizes add to the algebraic multiplicity $4$; the largest is the exponent $2$ in $m_T$.}
Sizes with largest $2$ and sum $4$: $\{2,2\}$ or $\{2,1,1\}$. So $J_2(2)\oplus J_2(2)$ or $J_2(2)\oplus J_1(2)\oplus J_1(2)$. They have the same $\chi_T$ and $m_T$; $\dim E(2,T)$ ($2$ versus $3$) tells them apart.
\end{prob}

\begin{prob}
\textbf{G5. }\par 
\idea{Over $\CC$, diagonalizable means $E(\lambda,T)=G(\lambda,T)$ for every $\lambda$.}
If $0\ne v\in G(\lambda,T)$, the hypothesis gives $Tv=\mu v$ for some $\mu$; then $0=(T-\lambda I)^nv=(\mu-\lambda)^nv$, so $\mu=\lambda$. Hence $G(\lambda,T)\subseteq E(\lambda,T)$, so $G(\lambda,T)=E(\lambda,T)$ for each eigenvalue. Then $V=\bigoplus G(\lambda_i,T)=\bigoplus E(\lambda_i,T)$ (Prop.~73), and a union of bases of the eigenspaces is an eigenbasis.
\end{prob}

\needspace{14\baselineskip}
\subsection{Trace (HW9)}

\begin{prob}
\textbf{HW9 P1. }\par 
\idea{Trace can be computed in any basis, so use the one where $P$ is simplest: $P$ is $0$ on $\nul P$ and the identity on $\range P$.}$V=\nul P\oplus\range P$. Take a basis $u_1,\dots,u_a$ of $\nul P$ and $w_1,\dots,w_b$ of $\range P$. $Pu_i=0$, and $w_j=Pv$ for some $v$ gives $Pw_j=P^2v=Pv=w_j$. So in the basis $u_1,\dots,u_a,w_1,\dots,w_b$, $\MM(P)=\operatorname{diag}(0,\dots,0,1,\dots,1)$ with $b$ ones, and $\tr P=b=\dim\range P=\operatorname{rank}P$.
\end{prob}

\begin{prob}
\textbf{HW9 P2. }\par 
\begin{enumerate}[(a)]
  \ii \idea{$m_T$ fixes the largest block for each eigenvalue and $\dim V=5$ fixes the total size; list the ways to fill $5$.}The largest block for $2$ has size $2$ and the largest block for $-1$ has size $1$. The eigenvalues of $T$ are exactly the roots $2,-1$ of $m_T$, so with $a=\dim G(2,T)\ge2$ and $b=\dim G(-1,T)\ge1$, Prop.~73 gives $a+b=5$: $(a,b)\in\{(2,3),(3,2),(4,1)\}$.
  \ii $\tr T=2a-b$: $1$, $4$, $7$.
  \ii $\tr=1$: $J_2(2)\oplus J_1(-1)^{\oplus3}$. $\tr=4$: $J_2(2)\oplus J_1(2)\oplus J_1(-1)^{\oplus2}$. $\tr=7$: $J_2(2)\oplus J_2(2)\oplus J_1(-1)$ (or $J_2(2)\oplus J_1(2)\oplus J_1(2)\oplus J_1(-1)$). Each has largest block $2$ for the eigenvalue $2$ and only $1\times1$ blocks for $-1$, so $m_T=(x-2)^2(x+1)$.
\end{enumerate}
\end{prob}

\begin{prob}
\textbf{HW9 P3. }\par 
\begin{enumerate}[(a)]
  \ii \idea{Trace is the sum of the eigenvalues, and a nilpotent operator has only the eigenvalue $0$.}$N$ is nilpotent, so its only eigenvalue is $0$ ($Nv=\lambda v$, $v\ne0$ gives $0=N^kv=\lambda^kv$). The trace is the sum of the eigenvalues with algebraic multiplicity, so $\tr N=0$. (Or: in a Jordan basis $\MM(N)$ has zero diagonal.)
  \ii $\operatorname{diag}(1,-1)$: trace $0$, and its square is $I$, so no power is $0$.
  \ii $A=\begin{psmallmatrix}0&1\\0&0\end{psmallmatrix}$, $B=\begin{psmallmatrix}0&0\\1&0\end{psmallmatrix}$: $AB=\begin{psmallmatrix}1&0\\0&0\end{psmallmatrix}$, $BA=\begin{psmallmatrix}0&0\\0&1\end{psmallmatrix}$, $AB-BA=\begin{psmallmatrix}1&0\\0&-1\end{psmallmatrix}\ne0$, with trace $1-1=0$.
\end{enumerate}
\end{prob}

\begin{prob}
\textbf{Tr1. }\par 
\idea{$A^2=-I$ pins down $m_A$; for a non-scalar $2\times2$ matrix, $m_A=\chi_A=x^2-(\tr A)x+\det A$, so compare coefficients.}
$m_A\mid x^2+1$. $A=cI$ would need $c^2=-1$ with $c$ real, impossible; so $\deg m_A=2$ and $m_A=x^2+1$. For a non-scalar $2\times2$ matrix, $m_A=\chi_A=x^2-(\tr A)x+\det A$. Comparing: $\tr A=0$, $\det A=1$.
\end{prob}

\begin{prob}
\textbf{Tr2. }\par 
\idea{The exponents in $m_T$ are the largest blocks; fill the dimension $3$.}
Algebraic multiplicities $a\ge1$ (for $1$) and $b\ge2$ (for $2$) with $a+b=3$: $a=1$, $b=2$. So $T\sim J_1(1)\oplus J_2(2)$, $\chi_T=(x-1)(x-2)^2$, $\tr T=1+2+2=5$.
\end{prob}

\needspace{14\baselineskip}
\subsection{Inner products (HW9)}

\begin{prob}
\textbf{HW9 P4. }\par 
\idea{Linearity and symmetry are usually immediate; test positivity by looking for a nonzero $v$ with $\ip vv\le0$.}$\ip xy_1$: yes. It is linear in $x$, symmetric, and $\ip xx_1=x_1^2+2x_2^2>0$ for $x\ne0$.
$\ip xy_2$: no. $\ip{e_2}{e_2}_2=-1<0$.
$\ip xy_3$: no. $\ip{(1,-1)}{(1,-1)}_3=0$ but $(1,-1)\ne0$. (It is bilinear and symmetric; only positivity fails.)
\end{prob}

\begin{prob}
\textbf{HW9 P5. }\par 
\begin{enumerate}[(a)]
  \ii \idea{Linearity and symmetry come from the integral; the content is positivity: a continuous function $\ge0$ with integral $0$ is $0$.}Linearity in $p$ comes from linearity of the integral; $\ip pq=\ip qp$; $\ip pp=\int_{-1}^1p^2\ge0$. If $\int_{-1}^1p^2=0$, the continuous function $p^2\ge0$ is $0$ on $[-1,1]$, so $p$ has infinitely many roots and $p=0$.
  \ii $\|1\|^2=2$, $\|x\|^2=\int x^2=\frac23$, $\|x^2\|^2=\int x^4=\frac25$: $\|1\|=\sqrt2$, $\|x\|=\sqrt{2/3}$, $\|x^2\|=\sqrt{2/5}$.
  \ii $\ip1x=\int x=0$ and $\ip x{x^2}=\int x^3=0$ (odd integrands), but $\ip1{x^2}=\frac23\ne0$. So $1\perp x$ and $x\perp x^2$; $1$ and $x^2$ are not orthogonal.
\end{enumerate}
\end{prob}

\begin{prob}
\textbf{HW9 P6. }\par 
\idea{Write the left side as an inner product with the all-ones vector, then apply Cauchy--Schwarz.}In $\RR^3$ with the dot product, $a+b+c=\ip{(a,b,c)}{(1,1,1)}$, so $(a+b+c)^2\le\|(a,b,c)\|^2\|(1,1,1)\|^2=3(a^2+b^2+c^2)$. Equality holds iff $(a,b,c)$ and $(1,1,1)$ are dependent, i.e.\ iff $a=b=c$.
\end{prob}

\begin{prob}
\textbf{HW9 P7. }\par 
\idea{Expand both squared norms; the cross terms appear with opposite signs.}$\|u+v\|^2=\ip uu+\ip uv+\ip vu+\ip vv$ and $\|u-v\|^2=\ip uu-\ip uv-\ip vu+\ip vv$. Adding, the cross terms cancel: $2\|u\|^2+2\|v\|^2$. This uses only (conjugate) bilinearity, so it holds over $\RR$ and $\CC$.
\end{prob}

\begin{prob}
\textbf{I1. }\par 
\idea{Find two vectors whose inner product is the constant $n$ and whose squared norms are the two factors: $u=(\sqrt{a_i})$ and $v=(1/\sqrt{a_i})$.}
In $\RR^n$, $\ip uv=\sum\sqrt{a_i}\cdot\frac1{\sqrt{a_i}}=n$, $\|u\|^2=\sum a_i$, $\|v\|^2=\sum\frac1{a_i}$. Cauchy--Schwarz: $n^2\le\big(\sum a_i\big)\big(\sum\frac1{a_i}\big)$. Equality iff $u=cv$, i.e.\ $\sqrt{a_i}=c/\sqrt{a_i}$, i.e.\ all $a_i$ equal $c$.
\end{prob}

\needspace{14\baselineskip}
\subsection{Orthogonality (HW10)}

\begin{prob}
\textbf{HW10 P1. }\par 
\idea{Subtract shadows, then normalize. On $[-1,1]$ odd integrands vanish, which kills many of the inner products.}$\|1\|^2=2$, so $e_1=\frac1{\sqrt2}$. $x-\ip x{e_1}e_1=x-\frac12\int_{-1}^1x\,dx=x$, and $\|x\|=\sqrt{2/3}$, so $e_2=\sqrt{\tfrac32}\,x$. $x^2-\ip{x^2}{e_1}e_1-\ip{x^2}{e_2}e_2=x^2-\frac12\cdot\frac23-0=x^2-\frac13$, with $\|x^2-\frac13\|^2=\int_{-1}^1\big(x^4-\frac23x^2+\frac19\big)dx=\frac25-\frac49+\frac29=\frac8{45}$. So $e_3=\sqrt{\tfrac{45}8}\big(x^2-\frac13\big)=\frac{\sqrt{10}}4(3x^2-1)$. Orthonormal basis: $\frac1{\sqrt2},\ \sqrt{\frac32}\,x,\ \frac{\sqrt{10}}4(3x^2-1)$.
\end{prob}

\begin{prob}
\textbf{HW10 P2. }\par 
\idea{Perpendicular to a span means perpendicular to each spanning vector: a homogeneous linear system.}$x\in W^\perp\iff x_1-x_2+x_3=0$ and $2x_2+x_3=0$. Then $x_3=-2x_2$, $x_1=x_2-x_3=3x_2$, and $x_4$ is free. Basis: $(3,1,-2,0)$, $(0,0,0,1)$. ($\dim W^\perp=4-2$.)
\end{prob}

\begin{prob}
\textbf{HW10 P3. }\par 
\idea{Same as Problem 2: impose $\ip p{1+x}=0$ on a general $p$.}For $p=a+bx+cx^2$: $\ip p{1+x}=\int_{-1}^1\big(a+(a+b)x+(b+c)x^2+cx^3\big)dx=2a+\frac23(b+c)$. So $p\in W^\perp\iff3a+b+c=0$. Basis: $1-3x$ ($a=1,c=0$) and $x^2-x$ ($a=0,c=1$). Check: $(1-3x)(1+x)=1-2x-3x^2$ integrates to $2-2=0$; $(x^2-x)(1+x)=x^3-x$ integrates to $0$. ($\dim W^\perp=3-1=2$.)
\end{prob}

\begin{prob}
\textbf{HW10 P4. }\par 
\idea{Prove the easier identity for sums, then take complements of both sides; $(Z^\perp)^\perp=Z$ undoes each $\perp$.}First, $(X+Y)^\perp=X^\perp\cap Y^\perp$ for any subspaces: $v\perp X+Y$ iff $v\perp X$ and $v\perp Y$ (every element of $X+Y$ is $x+y$, and $X,Y\subseteq X+Y$). Apply this to $X=U^\perp$, $Y=W^\perp$ and use $(Z^\perp)^\perp=Z$ (Prop.~90): $(U^\perp+W^\perp)^\perp=U\cap W$. Taking $\perp$ of both sides and using Prop.~90 again: $U^\perp+W^\perp=(U\cap W)^\perp$.
\end{prob}

\begin{prob}
\textbf{HW10 P5. }\par 
\begin{enumerate}[(a)]
  \ii \idea{The $U^\perp$ part of $v$ is invisible to inner products with the $e_j$.}Write $v=u+w$ with $u\in U$, $w\in U^\perp$ (Prop.~89); $P_Uv=u$. Since $w\perp e_j$, $\ip v{e_j}=\ip u{e_j}$. As $e_1,\dots,e_m$ is an orthonormal basis of $U$, $u=\sum_j\ip u{e_j}e_j=\sum_j\ip v{e_j}e_j$.
  \ii \idea{Split the error into a part in $U^\perp$ that every $u$ must keep and a part in $U$ you can remove; Pythagoras adds them in squares.}For $u\in U$, $v-u=(v-P_Uv)+(P_Uv-u)$ with $v-P_Uv\in U^\perp$ and $P_Uv-u\in U$. By Pythagoras $\|v-u\|^2=\|v-P_Uv\|^2+\|P_Uv-u\|^2\ge\|v-P_Uv\|^2$. Equality iff $\|P_Uv-u\|=0$, i.e.\ iff $u=P_Uv$.
\end{enumerate}
\end{prob}

\begin{prob}
\textbf{HW10 P6. }\par 
\idea{Riesz: find the representing vector by testing on an orthonormal basis; its $i$-th coordinate is $\psi(e_i)$ (no conjugates over $\RR$).}$\psi(f)=\int_0^2f$ is a linear functional, so Riesz gives a unique $p$. Orthonormal basis of $\PP_1(\RR)$: $e_1=\frac1{\sqrt2}$, $e_2=\sqrt{\frac32}x$ (HW10 P1). $p=\psi(e_1)e_1+\psi(e_2)e_2$ (real case, no conjugates). $\psi(e_1)=\frac2{\sqrt2}=\sqrt2$ and $\psi(e_2)=\sqrt{\frac32}\cdot2$. So $p=\sqrt2\cdot\frac1{\sqrt2}+2\cdot\frac32x=1+3x$. Check: $\ip1{1+3x}=2=\int_0^21$ and $\ip x{1+3x}=\int_{-1}^13x^2=2=\int_0^2x$.
\end{prob}

\begin{prob}
\textbf{HW10 P7. }\par 
\idea{($\Leftarrow$) If some $u\in\range P$ and $w\in\nul P$ were not perpendicular, moving $u$ slightly along $w$ would shorten it while $P$ still sends it to $u$, so $P$ would increase a length. To make this a proof, use $\|u+tw\|^2\ge\|u\|^2$ for all $t$ and choose $t$ pointing against $\ip uw$.}($\Rightarrow$) If $P=P_U$, then $v=P_Uv+(v-P_Uv)$ with orthogonal pieces, so $\|v\|^2=\|P_Uv\|^2+\|v-P_Uv\|^2\ge\|P_Uv\|^2$.
($\Leftarrow$) $V=\nul P\oplus\range P$ (Prop.~91). Let $u\in\range P$, $w\in\nul P$, and $t\in\FF$. Then $P(u+tw)=u$ (as $Pu=u$, $Pw=0$), so $\|u\|^2\le\|u+tw\|^2=\|u\|^2+2\operatorname{Re}\big(\bar t\ip uw\big)+|t|^2\|w\|^2$. Take $t=-s\ip uw$ with $s>0$ real: $0\le-2s|\ip uw|^2+s^2|\ip uw|^2\|w\|^2$. Divide by $s$ and let $s\to0^+$: $|\ip uw|^2\le0$, so $\ip uw=0$. Hence $\nul P\subseteq(\range P)^\perp$; both have dimension $\dim V-\dim\range P$, so $\nul P=(\range P)^\perp$. So for $v=u+w$ with $u\in\range P$, $w\in\nul P=(\range P)^\perp$: $Pv=Pu+Pw=u=P_{\range P}v$; hence $P=P_{\range P}$.
\end{prob}

\begin{prob}
\textbf{O1. }\par 
\idea{The integral is $\|x^2-p\|^2$ for the $L^2$ inner product on $[0,1]$, so the minimizer is the orthogonal projection of $x^2$ onto $\PP_1(\RR)$ (best approximation).}
Orthonormal basis of $\PP_1$ (Ex.~70): $e_1=1$, $e_2=\sqrt3(2x-1)$. $\ip{x^2}{e_1}=\frac13$, $\ip{x^2}{e_2}=\sqrt3\int_0^1x^2(2x-1)\,dx=\frac{\sqrt3}6$. So $p=\frac13+\frac{\sqrt3}6\cdot\sqrt3(2x-1)=x-\frac16$. Minimum: $\int_0^1\big(x^2-x+\frac16\big)^2dx=\frac1{180}$ (this is $\|w_3\|^2$ in Ex.~70).
\end{prob}

\begin{prob}
\textbf{O2. }\par 
\idea{Projection onto a line: $\frac{\ip vu}{\|u\|^2}u$. $U^\perp$ is one linear equation.}
$P_Uv=\frac{10}4(1,1,1,1)=\big(\frac52,\frac52,\frac52,\frac52\big)$. $U^\perp=\{x_1+x_2+x_3+x_4=0\}$, basis $(1,-1,0,0),(0,1,-1,0),(0,0,1,-1)$. Distance $\|v-P_Uv\|=\big\|\big(-\frac32,-\frac12,\frac12,\frac32\big)\big\|=\sqrt5$.
\end{prob}

\begin{prob}
\textbf{O3. }\par 
\idea{Split $v=u+w$ along $U\oplus U^\perp$: $R$ keeps $w$ and flips $u$, a reflection.}
$Rv=v-2P_Uv=(u+w)-2u=-u+w$, so $\|Rv\|^2=\|u\|^2+\|w\|^2=\|v\|^2$ (Pythagoras). $R^2=I-4P_U+4P_U^2=I$ since $P_U^2=P_U$.
\end{prob}

\end{document}
